% Scheletro di lavoro, non testo definitivo.
% Fonti e passi sono proposte di lettura: verificarli sul testo greco prima della stesura.
% Seguira il metodo: documento originale, traduzione, esame dell'argomentazione, ricostruzione critica.
\chapter{La teoria delle idee}
\label{cap:platone:idee}

% Problema: Quali difficoltà della conoscenza e del divenire motivano la distinzione?
% Fonti da verificare: Fedone 74a–76e, 99d–102a; Repubblica V 475e–480a.
\section{Perché postulare le forme?}

\subsection{L'uguale e l'Uguale in sé}

\subsection{Le cose molteplici e ciò che esse sono}

% Problema: Che cosa significa che una cosa partecipa di un'idea?
% Fonti da verificare: Fedone 100b–102a; Repubblica V 476a–480a.
\section{Il rapporto fra sensibili e intelligibili}

\subsection{La partecipazione e la predicazione}

\subsection{Unità della forma e molteplicità delle cose}

% Problema: Quali obiezioni Platone mette in scena alla separazione delle forme?
% Fonti da verificare: Parmenide 130a–135d.
\section{Le aporie della teoria}

\subsection{L'estensione delle idee e il problema della partecipazione}

\subsection{Regresso, conoscibilità e difficoltà della separazione}

% Problema: Quale ruolo ha il Bene rispetto alle altre idee?
% Fonti da verificare: Repubblica VI 504a–509c; raccordo con cap. 26 e 29.
\section{L'idea del Bene}

\subsection{Il Bene come problema}

\subsection{Limiti della ricostruzione sistematica}
